\chapter{Ход работы}
\label{ch:intro}

\section{Архитектура проекта}

\begin{figure}[H]
    \centering
    \includegraphics{ARCH.png}
    \caption{Архитектура проекта}
\end{figure}

Компоненты приложения API и DATABASE находятся в разных контейнерах и общаются между собой. Оба контейнера находятся в POD (минимальная абстрактная единица в kubernetes).
В одном kubernetes кластере может быть несколько подов, нагрузка между ними может балансироваться. Всё это развернуто на сервере. Пользователь может подклчюиться к API
и сделать запрос на получение данных из базы данных (GET), запрос на создание новой записи в базе данных (POST), обновить запись (PUT). Разработчик может вносить изменения в проект,
делать push в удаленный репозиторий, где запускается CI/CD pipeline, в котором происходит сборка docker образа и деплой на сервер.

\section{REST}

REST API (Representational State Transfer Application Programming Interface) — это архитектурный стиль взаимодействия между клиентом и сервером через 
HTTP. Он определяет принципы построения API, обеспечивая стандартизированный и эффективный обмен данными между различными системами.

REST API строится на нескольких ключевых принципах:
\begin{itemize}
    \item \textbf{Клиент-серверная архитектура.}  
    REST API подразумевает чёткое разделение между клиентом и сервером. Клиент запрашивает данные, а сервер их предоставляет. Такое разделение улучшает масштабируемость системы и позволяет клиентам работать независимо от серверной логики.

    \item \textbf{Отсутствие состояния (Stateless).}  
    Каждый запрос от клиента к серверу обрабатывается независимо от предыдущих запросов. Сервер не хранит информацию о состоянии клиента, что упрощает обработку запросов и повышает отказоустойчивость системы.

    \item \textbf{Кеширование.}  
    Ответы сервера могут кешироваться, чтобы снизить нагрузку на сервер и ускорить загрузку данных. REST API поддерживает механизмы кеширования, такие как HTTP-заголовки \texttt{Cache-Control} и \texttt{ETag}, которые позволяют клиентам повторно использовать ранее полученные данные.

    \item \textbf{Единообразие интерфейса.}  
    Все ресурсы REST API должны иметь чёткую структуру и единообразные URL-адреса. Запросы к API выполняются с использованием стандартных методов HTTP, а данные передаются в предсказуемых форматах, таких как JSON или XML.

    \item \textbf{Система уровней (Layered System).}  
    REST API может включать несколько уровней, таких как балансировщики нагрузки, прокси-серверы и системы аутентификации. Каждый уровень выполняет свою функцию и не зависит от других, что повышает надёжность и гибкость системы.

    \item \textbf{Возможность выполнения кода по требованию.}  
    Хотя этот принцип не является обязательным, REST API может поддерживать загрузку и выполнение кода на стороне клиента, например в виде скриптов или небольших программ.
\end{itemize}

\section{Разработанное API}

\subsection{Общее описание}
Разработанное API представляет собой посредника между пользователем и базой данных: пользователь хочет получить/загрузить данные, отправляет запрос к API, которое
проверяет запрос, переводит на язык запросов SQL, отправляет SQL запрос в базу данных, получает данные (или вносит) и отправляет ответ пользователю. Такая
схема позволяет пользователю не задумываться о каком-то там языке запрос SQL, а просто сказать API: "Дай мне вот это". \\

\subsection{Техническое описание}
API реализовано на языке Python с помощью библиотеки fastapi и psycopg (для работы с базой данных). \\

API работает по следующему принципу:
\begin{enumerate}
    \item Пытается подключиться к БД, если попытка неудачная - попытаться подключиться снова.
    \item Создание endpoint'ов и ожидание запросов на них.
\end{enumerate}

\subsubsection{Описание endpoint'ов}

\begin{enumerate}
    \item GET("/books") - возвращает все книги из базы данных
    \item GET("/books/{id}") - возвращает определенную книгу по ID
    \item POST("/books") - делает новую запись в БД
    \item PUT("/books/{id}") - обновление записи в БД
    \item DELETE("/books/{id}") - удаление книги по ID
\end{enumerate}

\subsection{База данных}
К базе данных нет каких-то тонких требований, поэтому я выбрал наиболее популярную и легковесную версию: postgress:alpine.

\section{Docker}
Docker - платформа для контейнеризации, которая позволяет упаковать приложение со всем его окружением (библиотеками, зависимостями, настройками) в 
изолированный контейнер. Это позволяет запускать много разных приложений без возникновнения проблем с зависимостями. \\

В докере собирается только код на Python (без БД). В качестве базового образа для сборки используется python:slim, который уже имеет предустановленный python и весит немного.
Установка зависимостей выполняется с помощью файла requirements.txt, который содержит все необходимые зависимости.

\subsection{Docker compose}
Docker compose позволяет поднимать несколько контейнеров разом. Это быстрее и удобнее, чем отдельно поднимать каждый контейнер.
Создание Docker compose позволяет локально собрать API вместе с БД для проверки.

\section{Локальный запуск}
Поднимем приложение и попытаемся залить данные в базу данных:

\begin{figure}[H]
    \centering
    \includegraphics{postman.png}
    \caption{Демонстрация локального запуска}
\end{figure}

Все работает успешно, поэтому можно переходить к kubernetes.

\section{Подготовка манифестов kubernetes}

Kubernetes (K8s) - открытая платформа для автоматизации развертывания, масштабирования и управления контейнеризированными приложениями. \\

Конфигурация kubernetes выполняется в специальных файлах, называемых манифестом, написанных в .yaml форматах. При конфигурировании используется декларативный стиль, т.е нет необходимости
вручную прописывать каждое действие, требуется только написать, что мы хотим сделать, а реализацию этого всего берет на себя kubernetes. \\

Для минимальной конфигурации необходимо минимум 3 файла: deployment.yaml, ingress.yaml, service.yaml. \\

\begin{enumerate}
    \item deployment.yaml cоздает и обновляет указанное количество копий контейнеров (Pods). Здесь описываются контейнеры (образы, кол-во реплик), окружения и тд.
    \item service.yaml предоставляет постоянный IP-адрес и DNS-имя для группы подов. В этом есть необходимость, поскольку поды могут пересоздаваться и их IP адрес будет меняться.
    Это значит, что коммуникация между подами может нарушиться.
    \item ingress.yaml принимает внешний HTTP/HTTPS трафик и направляет его внутрь кластера,
\end{enumerate} 

\begin{figure}[H]
    \centering
    \includegraphics{kuber.png}
    \caption{Взаимосвязь сущностей kubernetes между собой}
\end{figure}

\subsection{Подготовка helm chart}
Helm - пакетный менеджер для Kubernetes. Это как apt в ubuntu, но только для kubernetes. Пакетный менеджер позволяет быстро и просто установить 
пакет (chart) и развернуть на сервере без необходимости писать манифесты (они вшиты в char). Помимо этого helm предоставляет функционал для управления, контроля релизов 
установленных пакетов, что очень удобно. \\

Chart - это пакет для helm. Т.е мы можем определенным образом упаковать наши манифесты в chart, а другой человек сможет с помощью helm удобно этот chart
развернуть на сервере. \\

Для создания chart необходимо сделать такую структуру директорий:

\begin{figure}[H]
    \centering
    \includegraphics{chart.png}
    \caption{Структура chart}
\end{figure}

Здесь файл chart.yaml отвечает за хранение служебной информации (по сути просто говорит, что здесь находится chart, чтобы helm его увидел). В файле values.yaml
определены значения для переменных для манифестов, которые можно менять и которые будут подставляться в манифесты автоматически. Таким образом, если мы 
хотим изменить какие то параметры, то не нужно лезть в манифесты, а достаточно поменять только файл values.yaml. В директории templates собраны все манифесты.

\section{CI/CD}
CI/CD (Continuous Integration/Continuous Delivery) - набор практик автоматизации, ускоряющий доставку кода от разработчика до продакшена. 
CI обеспечивает автоматическую сборку и тестирование, а CD — надежный выпуск продукта, минимизируя ручные операции, устраняя ошибки и повышая стабильность разработки. \\

Pipeline CD/CD состоит из следующих шагов:
\begin{enumerate}
    \item Сборка docker образа и push в ghcr.io (github хранилище docker образов. Отсюда при деплое будет подтягиваться последняя версия приложения)
    \item Деплой на сервер
\end{enumerate}

Pipeline начинает выполняться только в случае, если был сделан push с тегом. Это сделано, чтобы CI/CD не триггерился на каждый коммит, а только на тегированные (предусмотренные).

\section{Недостатки работы}
\begin{enumerate}
    \item Никакой безопасности.
    \item Отсутствие Persistent Volume (БД очищается после пересоздания пода)
\end{enumerate}
\endinput